% !TeX root = ../papers/nc-p3.tex
\section{階数一四次成分の滑らかさと楕円四次の二次式鉛筆}
\label{sec:rank-one-quartic-normalization}
\begin{theorem}\label{thm:rank-one-quartic-normalization}
被約窓の一般可逆な全階数反対根の整な階数一四次成分は、存在するなら滑らかな楕円正規四次である。
二埋込みの超平面束の差は非自明な二捩れである。曲線外の印は少なくとも八つあり、
二次式鉛筆の一つの曲面上では二印まで、例外的にその母線全体が階数一の場合に三印までである。
この例外の鉛筆元は高々一つである。
\end{theorem}
\begin{proof}[四次の二つの座標表示]
平面曲線は二次までなので四次成分 $C$ は非退化で、算術種数はCastelnuovo上界1以下である。
次数四の線束の二乗の切断数は正規化上で9以下なので二次式のイデアルは非零である。
独立な二次式が二つあれば既約な二曲面の完全交叉を $C$ が次数で使い切り、算術種数1の完全交叉そのものとなる。
この場合 $h^0(\mathcal O_C(1))=4$ で完全な埋込みである。逆に算術種数1なら
$h^0(\mathcal O_C(2))=8$ なので二つの二次式がある。両側で二次評価の核次元は等しい。

二つの埋込みが射影的に一致する場合は不可能である。実際 $N=xx^t+A(q)$ で、
$A(q)$ は高々二つの一定対称行列による鉛筆である。一般階数4なら階数二点は高々二つで、
$x\wedge y$ の六つの独立方向を与えられない。一般階数3なら余因子の二項三次式の原始核は一定か一次である。
固定核なら $Nv=(x^tv)x$、動く場合はスカラー因子の次数1ゆえやはり六方向はない。
一般階数2以下では $N$ が一般可逆でなくなる。この議論は係数空間が一次元以下の場合も含む。

二次式の空間が一つの場合、正規化は有理である。いずれの埋込みの座標も正規化上の
$H^0(\mathbb P^1,\mathcal O(4))$ の四空間で、少なくとも三つの独立な共通切断がある。
この三切断を両側の最初の座標に取り、対応する三行三列小行列を $N_0$、入力三座標を $b$ とすると
$N_0=bb^t+qA$、$A$ は三行定数行列である。$q$ の外の印で階数差を読むと $A$ は階数2以下。
$0\ne w\in\ker A$ と置く。印で $N_0=cc^t$ と自明化すると $(c^t w)c=(b^tw)b$ かつ
$(c^tw)^2=(b^tw)^2$ であり、$b^tw\ne0$ の印では $cc^t=bb^t$ になる。
$A\ne0$ なら $q(b^tw)$ が全印で消えて三次評価に反する。$A=0$ なら共通二座標の交代積が全印で消えて
混合評価に反する。これで一つの二次式の場合を除いた。
\end{proof}
\begin{proof}[曲線外の印の鉛筆パラメータ]
完全交叉の場合 $L'\ne L$ が必要で、$\eta=L'/L$ は非自明な二捩れである。
$h^0(LL')=8$ なので混合評価から曲線外の印が八つ以上必要である。
$C$ の二乗の線型包は二つの二次式の零化空間で、共通の十次元座標の商二空間が両鉛筆の
パラメータを同定する。各曲線外の点には、それを含む唯一の鉛筆元のパラメータを割り当てる。

同じ $Q_t$ 上の異なる二つの曲線外印を出力で結んだ直線が $Q_t$ に入らなければ $C$ と交わらない。
実際二つの印が既に直線と二次曲面の交数を使い切る。この直線を通る一般平面は四つの動点を切り、
平面鉛筆への射の一般繊維なのでそれらは一つの体の共役点である。
四動点と二印は $Q_t$ の平面二次曲線上にあるから六平方に関係がある。二印のみの関係はなく四動点は全て台に入る。
全等方性から入力四点も共面となり、抽象曲線上で $L=L'$ になる。一般平面は特異点を避けるためこの議論には
滑らかさを要しない。従って同じ群のどの二印を結ぶ直線も $Q_t$ に含まれる。
全鉛筆元は既約であり、三つ以上の非共線な点の三辺を同じ既約二次曲面は含めない。
よって各群は共線で高々三印。三印ならその母線はもう一つの二次式で $C$ を長さ二で切る。
五つ分の消滅から逆向きの四次小行列式は母線全体で零となる。総次数上界より
そのような追加直線は高々一本で、直線は鉛筆のただ一元に含まれるため三印の群も高々一つである。
\end{proof}
\begin{proof}[有理正規化を持つ完全交叉]
完全交叉が特異なら算術種数1なので正規化は有理で、先と同じ三つの共通切断を取る。
今は $N_0=bb^t+A(q_1,q_2)$ で三行の定数対称行列鉛筆となる。
一般可逆なら $C$ 外の各印でこの三行鉛筆の階数は2以下なので、曲線外の各群のパラメータは三つの根に限られる。
前段の群上界は合わせて $3+2+2=7$ で不足する。一般に固定核 $w$ があれば、先の三座標の議論により
どの非零の補正成分 $q_0$ についても $q_0(b^tw)$ が全印で消え、補正が零なら交代積が消えるので不可能となる。

固定核のない三行特異鉛筆の一般階数は2である（階数一対称行列の射影直線はない）。
二次余因子の原始核は一次で基点がなく、二項変数を $s,t$ として基底で $(s,t,0)^t$ と書ける。
一次の各行の最初の二成分は $(t,-s)$ の定数倍、対称性より先頭二行二列は零である。
従ってその二成分では $N_0$ が $bb^t$ に厳密に一致し、全印で非零の混合交代式が消えてしまう。
これで特異完全交叉を排除した。残る算術種数1の滑らかな場合が主張の楕円四次である。
\end{proof}
